% EM structure and teaching revision 20260928
\subsection{回旋加速的同步与失谐}
理想回旋加速器的匀强磁场只负责弯转，交变电场在电极缝隙中做功。非相对论粒子每半周回到缝隙，时间为 $\pi m/(|q|B)$，所以高频电压必须在这段时间内反号，驱动频率等于回旋频率。
若每次过缝时电压大小近似为 $U_g$，且粒子总在加速相位通过，则每次获得动能 $|q|U_g$。经过 $N$ 次过缝，忽略初能与损失，有
\[
 K_N=N|q|U_g,\qquad r_N=\sqrt{\frac{2mNU_g}{|q|B^2}}.
\]
因而半径随 $\sqrt N$ 增长，轨道间距逐渐变小。半径达到装置外缘 $R$ 时，经典最大动能为 $q^2B^2R^2/(2m)$，它同时受到装置尺寸与磁场强度限制。
高能区中 $\gamma$ 增大，回旋周期随之变长，恒频电源与粒子到缝时间失配。此时即使磁场强度和电极电压不变，粒子也可能在不合适的相位减速；调整驱动频率或磁场才可能恢复同步。

